第一编 8.251 需要微积分、基础力学和量子谐振子, 并在正文补足相对论、变分和自由场知识. 第二编 8.821 还需要量子场论、规范场与广义相对论基础. 曲率、黑洞因果结构承接本仓库广义相对论卷; 规范场、路径积分及重整化承接量子场论卷. 本卷只在弦和全息计算需要的地方解释相应工具. 约定目标时空度规号差为 $(-,+,\ldots,+)$; 除明确恢复单位处, 使用 $\hbar=c=1$.

\begin{itemize}
\item 8.251 String Theory for Undergraduates, Spring 2007, Barton Zwiebach、Alan Guth. 26 讲笔记由 Bryan Owens 根据 Kayla Jacobs 的课堂手稿编成, 官方明确说明 Zwiebach 没有全面校阅. 本卷重算作用量、约束、谱与相应题解, 不把学生笔记的公式直接视为无误. 十一份作业中大量项目只给 2004 年 Cambridge 教材题号; 这些题面不猜补. 公开作业中确实完整重述且无相反许可标记的题面逐件检查后纳入, 保留其教材引用与教师署名. 仅有教材题号的项目仍只留定位; 不取用第三方教材原本补造缺失条件.
\item 8.821 String Theory and Holographic Duality, OCW 页面标 Fall 2014, Hong Liu. 页面列 24 讲、五组作业, 前四组有官方解答, 无考试. 原文件页脚有 Spring 2015 字样, 作业正文又标 Fall 2014; 两种标记保留在来源清单, 不擅自统一. 本卷按页面课程版本署名. 前四组题解核对官方答案, 第五组题解由 AI 独立推导.
\end{itemize}

正文为中文主题讲义与独立推导, 不是原文件逐页翻译. 当前正文逐讲对应两门课的 50 个讲次, 包括第 26 讲三道复习题. 所有必做 8.821 作业已有完整题面与解答草稿. 8.251 三份正式试卷、两份历史试卷及完整公开作业题均已写入草稿, 未取得教材题面仍逐项列明. 本稿的正文例子为自编说明, 不冒充官方习题或计入其题数. 逐讲逐题稿尚待内容与排版复核, 本稿不能作为已验收的完整交付版本. 世界面经典理论、自由玻色弦谱和经典全息计算是主线. 超弦只展开自由谱与 GSO 投影所需知识, 不声称完成超弦相互作用、模空间积分、超引力或非微扰弦论. AdS/CFT 作为对偶猜想使用; 经典引力近似中的计算是对猜想的后果检验, 不是猜想的一般证明.

基于 MIT OCW 的改编及独立增补采用 \href{https://creativecommons.org/licenses/by-nc-sa/4.0/}{CC BY-NC-SA 4.0}, 须署名、非商业、相同方式共享. 中文翻译、重排、勘误说明和自绘图不表示 MIT 或作者认可. 不复制第三方书本、论文全文与受限图像; 逐文件来源、许可例外与未覆盖范围另见源码包中的 \texttt{sources/} 和 \texttt{coverage.md}.
